%
% Halaman Abstract
%
% @author  Andreas Febrian
% @version 2.2.0
% @edit by Ichlasul Affan
%

\chapter*{ABSTRACT}
\singlespacing

\noindent \begin{tabular}{l l p{11.0cm}}
	\ifx\blank\npmDua
		Name&: & \penulisSatu \\
		Study Program&: & \studyProgramSatu \\
	\else
		Writer 1 / Study Program&: & \penulisSatu~/ \studyProgramSatu\\
		Writer 2 / Study Program&: & \penulisDua~/ \studyProgramDua\\
	\fi
	\ifx\blank\npmTiga\else
		Writer 3 / Study Program&: & \penulisTiga~/ \studyProgramTiga\\
	\fi
	Title&: & \judulInggris \\
	Counselor&: & \pembimbingSatu \\
	\ifx\blank\pembimbingDua
	\else
		\ &\ & \pembimbingDua \\
	\fi
	\ifx\blank\pembimbingTiga
	\else
		\ &\ & \pembimbingTiga \\
	\fi
\end{tabular} \\

\vspace*{0.5cm}

Dynamic Kubernetes clusters accumulate multi-dimensional resource fragmentation that strands capacity, while compute-focused consolidation may disrupt busy workloads or move dependent microservices across degraded network links. This study presents a Kubernetes Descheduler architecture that combines resource defragmentation with actual-usage and network-aware eviction safeguards. The proposed \textit{ResourceDefragmentation} policy evaluates CPU-memory shape using the Resource Imbalance Index (RII), Free-Space Index (FSI), and bin scores, then selects evictions based on feasible predicted scheduler targets. Actual-usage and network-aware filters respectively protect busy Pods and block movements toward latency-degrading destinations. Across five request-based resource scenarios, the policy reclaimed four Nodes in the underutilized scenario and three Nodes in the fragmented scenario. In the heterogeneous-drain scenario, it reduced Mean Absolute RII from the low-20\% range to around 2\% using two evictions. In the actual-usage evaluation, the filter blocked eviction of a busy API Pod and redirected the strategy to an idle candidate, reducing the application failure rate from \(6.57\%\) to \(0.00\%\). In the network-aware evaluation, the baseline ResourceDefragmentation behavior balanced compute resources but evicted sensitive workloads, causing latency spikes of up to $\sim$167ms. In contrast, the latency-aware filter intercepted these network-degrading actions, maintained baseline latency at $\sim$8ms, and still reduced Mean Absolute RII from the 50\% range to around 17\% by steering eviction toward non-critical Pods. Overall, the proposed architecture improves resource utilization while protecting active and network-sensitive workloads.

\vspace*{0.2cm}

\noindent Key words: \\ Kubernetes, Descheduler, Resource Defragmentation, Runtime Resource Usage, Network Latency, Microservices \\

\setstretch{1.4}
\newpage
